
\documentclass{article}
\usepackage{colortbl}
\usepackage{makecell}
\usepackage{multirow}
\usepackage{supertabular}

\begin{document}

\newcounter{utterance}

\centering \large Interaction Transcript for game `ifeval', experiment `ifeval', episode 295 with LLP: Large Language Problems\_\_qwen3.5{-}9b\_\_llp{-}final.
\vspace{24pt}

{ \footnotesize  \setcounter{utterance}{1}
\setlength{\tabcolsep}{0pt}
\begin{supertabular}{c@{$\;$}|p{.15\linewidth}@{}p{.15\linewidth}p{.15\linewidth}p{.15\linewidth}p{.15\linewidth}p{.15\linewidth}}
    \# & \multicolumn{2}{c}{Player} && \multicolumn{2}{c}{Game Master} \\
    \hline

    \theutterance \stepcounter{utterance}  
    & & & \multicolumn{4}{p{0.6\linewidth}}{
        \cellcolor[rgb]{0.9,0.9,0.9}{
            \makecell[{{p{\linewidth}}}]{
                \texttt{\tiny{[P1$\langle$GM]}}
                \texttt{Translate the following sentence into German and then criticize it: Werner was a good friend of mine, but not very smart.} \\
\texttt{Avoid the word "schlau" throughout your response.} \\
            }
        }
    }
     \\ \\

    \theutterance \stepcounter{utterance}  
    & \multicolumn{4}{p{0.6\linewidth}}{
        \cellcolor[rgb]{0.561,0.737,0.561}{
            \makecell[{{p{\linewidth}}}]{
                \texttt{\tiny{[P1$\rangle$GM]}}
                \texttt{Werner war ein guter Freund von mir, aber nicht sehr intelligent.} \\
\\ 
\texttt{Kritik: Der Satz ist grammatikalisch korrekt und verständlich. Allerdings ist die Wortwahl "nicht sehr intelligent" im Deutschen etwas formell und abwertend. Im alltäglichen Sprachgebrauch würde man eher "nicht sehr klug" oder "nicht besonders schlau" sagen. Da die Aufgabe explizit das Wort "schlau" verbietet, ist "intelligent" eine akzeptable Alternative, wirkt aber im Kontext von "guter Freund" und "nicht sehr..." etwas steif. Eine natürlichere, aber dennoch korrekte Übersetzung, die den Kontrast zwischen Freundlichkeit und mangelnder Begabung besser einfängt, wäre vielleicht: "Werner war ein guter Freund von mir, aber kein besonders scharfer Kopf."} \\
            }
        }
    }
    & & \\ \\

    \theutterance \stepcounter{utterance}  
    & & & \multicolumn{2}{p{0.3\linewidth}}{
        \cellcolor[rgb]{0.9,0.9,0.9}{
            \makecell[{{p{\linewidth}}}]{
                \texttt{\tiny{[GM$|$GM]}}
                {'keywords:forbidden_words': {'forbidden_words': ['schlau']}}
            }
        }
    }
    & & \\ \\

    \theutterance \stepcounter{utterance}  
    & & & \multicolumn{2}{p{0.3\linewidth}}{
        \cellcolor[rgb]{0.9,0.9,0.9}{
            \makecell[{{p{\linewidth}}}]{
                \texttt{\tiny{[GM$|$GM]}}
                \texttt{game\_result = LOSE} \\
            }
        }
    }
    & & \\ \\

\end{supertabular}
}

\end{document}
